\newpage
\section{Related Work}
\label{ch:related-work}

\subsection{Limitations of the Base Study}
\label{sec:rw-base-study}

While Go et al.~\cite{go2019} demonstrated the feasibility of an automated system for identifying and classifying viable pollen grains, their proposed image-analysis protocol encountered difficulties when pollen grains overlapped or occurred in clustered (``lumped'') arrangements. The method used edge detection, dilation, region filling, and $k$-means clustering, and relied on manually selected image-processing parameters, limiting its adaptability to variations in pollen images. In particular, accuracy decreased from 98.3\% for images without lumped pollen to 85.06\% for images containing lumped pollen because the edge-detection process failed to consistently distinguish overlapping grains. These limitations highlighted the need for more adaptive machine-learning-based methods. Subsequent studies have explored machine learning, deep learning, segmentation, and transfer-learning techniques to address these limitations, providing a basis for determining whether the identified gap has been addressed, partially addressed, or remains unresolved.

\subsection{Identification and Detection of Pollen Grains}
\label{sec:rw-detection}

Recent studies have increasingly relied on machine learning and deep learning to reduce the limitations of manually designed image-processing procedures. Rather than depending on predefined parameters, as in Go et al.~\cite{go2019}, these approaches attempt to learn useful pollen features directly from image data.

Viertel and König~\cite{viertel2022} reviewed automated pollen-image classification methods and compared traditional machine-learning approaches with convolutional neural networks (CNNs). Their review highlighted the increasing use of deep learning while also identifying continuing challenges involving image acquisition, dataset diversity, and model generalization. Chen and Ju~\cite{chen2022} addressed image-quality variation by combining SRGAN-based deblurring with a multi-scale classifier. Their ResNet-50-based model achieved 97.7\% accuracy, demonstrating that improving the quality of microscopic images before classification can improve recognition performance. Jin et al.~\cite{jin2023} further extended automated pollen detection to partially labeled datasets through Faster R-CNN and a teacher-student semi-supervised learning strategy. Their approach achieved an F1 score of 76.9\%, compared with 61.3\% for the BAA500 automated pollen monitor, showing that machine learning can improve detection even when complete manual annotation is unavailable. Shamrat et al.~\cite{shamrat2024} also developed PollenNet, which combined deep learning with preprocessing procedures such as removal of low-quality and duplicate images, noise reduction, brightness adjustment, and explainable artificial intelligence.

These studies demonstrate a progression from fixed image-processing parameters toward models that learn features from data and can accommodate different image conditions. However, image quality, incomplete labeling, and dataset diversity continue to influence detection performance.

\subsection{Detection and Segmentation of Overlapping or Clustered Pollen Grains}
\label{sec:rw-overlapping}

The difficulty of identifying lumped pollen grains is the most direct limitation of Go et al.~\cite{go2019}. Because dilation and region-filling operations depended on correctly detected pollen boundaries, failures during detection affected the final pollen count and viability classification. This limitation has encouraged researchers to use segmentation and object-detection approaches that can distinguish individual objects within complex images.

Li et al.~\cite{li2023weakly} proposed a weakly supervised collaborative-learning framework that combined pollen detection, U-Net segmentation, and CNN classification. The method generated segmentation information from pollen contours and allowed segmentation and classification to improve each other. Using approximately 1,000 images with image-level labels, the method achieved 86.6\% pollen-identification accuracy, demonstrating that segmentation can reduce interference from impurities while also reducing the need for extensive pixel-level annotation.

Zhang et al.~\cite{zhang2023} used Mask R-CNN instance segmentation to detect pollen germination and pollen-tube length from microscopic images. To improve recognition of difficult structures, synthetic images containing crossed pollen tubes were added to the training data. The model achieved an mAP of 0.949 at an IoU threshold of 0.50, while its measurements of germination frequency and pollen-tube length produced $R^2$ values of 0.909 and 0.958, respectively. The model was also successfully applied to two additional plant species, indicating some capacity for generalization.

Gallardo et al.~\cite{gallardo2024} approached the problem through multifocus microscopy combined with deep-learning object detection. Faster R-CNN and RetinaNet were used to process multiple focal planes, improving the detection of grains that might be blurred or obscured in individual images. The final system achieved a macro-F1 score of 0.956. However, highly overlapping grains and multiple detections across focal planes could complicate non-maximum suppression, showing that multifocus imaging can introduce its own detection challenges.

Feng et al.~\cite{feng2025} further demonstrated that pollen segmentation can be affected by differences between imaging domains. Their study investigated mixture-of-experts models, domain fine-tuning, continual learning, and foundation-model segmentation to detect and isolate pollen on palynological slides. Fine-tuning substantially improved detection when models were transferred between imaging domains, while segmentation models were able to further refine detected pollen boundaries. These findings show that improving segmentation alone may not be sufficient when a model is applied to substantially different datasets.

Overall, recent object-detection and segmentation techniques have substantially improved the handling of complex pollen arrangements compared with the edge-based method of Go et al.~\cite{go2019}. Nevertheless, overlapping grains, different imaging conditions, focal-plane variation, and insufficiently diverse training data remain challenges.

\subsection{Machine Learning for Adaptive Pollen Viability Classification}
\label{sec:rw-viability}

Go et al.~\cite{go2019} classified viable and nonviable pollen using $k$-means clustering based on grayscale intensity after staining. Although effective for relatively separated grains, the process depended on the accuracy of the preceding image-processing stages. The authors consequently suggested that machine learning could make pollen identification and processing parameters more dynamic.

Recent studies directly address this limitation through deep learning. Tan et al.~\cite{tan2022} developed PollenDetect, a modified YOLOv5 model designed to detect small pollen targets and distinguish viable from nonviable stained pollen. The system achieved 99\% accuracy while reducing processing time from approximately three minutes to about one second per image. It also performed better than YOLOv3 and Faster R-CNN and showed improved detection of occluded and overlapping pollen aggregates. However, the system required additional training and parameter adjustment when applied to other plant species, illustrating a remaining transferability limitation.

Yamazaki et al.~\cite{yamazaki2023} approached viability through pollen germination rather than staining. Their YOLOv5 transfer-learning model distinguished germinated from non-germinated pollen and estimated germination rates across several plant species. Together with the germination measurements of Zhang et al.~\cite{zhang2023}, this shows that machine learning can move beyond simple viable/nonviable classification and support more detailed quantitative measurements of pollen viability.

Other researchers have focused on making pollen classification more robust under difficult image conditions. Li et al.~\cite{li2024prior} developed a prior knowledge-guided deep feature-learning model that incorporated pollen location, shape, and texture information to guide CNN attention. The approach achieved accuracy above 88\% on real-world pollen images, demonstrating the value of combining domain-specific pollen characteristics with deep learning. Zolfaghari and Sajedi~\cite{zolfaghari2024} similarly developed RCANet, which uses separate attention mechanisms for pollen structure and spatial information. Across three pollen datasets, the model achieved weighted F1 scores of 98.69\%, 97.83\%, and 98.24\%.

The most recent work further targets challenging imaging conditions. Li et al.~\cite{li2026pvdt} developed PVDT, which combines YOLOv10-CBAM, pollen-image segmentation, and attention-based processing. The system specifically considered variable lighting, high pollen density, and low image quality, achieving 90--98\% accuracy under these conditions. Transfer learning also enabled pollen-viability detection across ten additional plant species.

These studies indicate that machine learning has significantly advanced the adaptability of automated pollen analysis. However, differences in species, pollen density, image quality, and imaging environments can still affect performance.

\subsection{Research Gap Analysis and Possible Research Direction}
\label{sec:rw-gap-analysis}

While subsequent studies have demonstrated that deep-learning object detection and instance segmentation can more reliably isolate complex and overlapping pollen grains~\cite{zhang2023, gallardo2024, tan2022, li2026pvdt}, these approaches have not been applied to or evaluated on rice pollen viability classification. Go et al.~\cite{go2019} remain the only reviewed study addressing automated viability identification of rice pollen, and their edge-based protocol lost accuracy when pollen grains overlapped or occurred in lumped arrangements. The deep-learning methods that better handle such arrangements were developed for other species or pollen types and frequently required additional training before being transferred to new crops~\cite{tan2022, feng2025}. Furthermore, apart from Go et al.~\cite{go2019}, the reviewed studies do not report viability performance separately for isolated and overlapping grains, leaving the effect of overlap on classification unquantified.

To address this gap, this study proposes to fine-tune a pretrained YOLO instance-segmentation model to detect individual rice pollen grains and classify each as viable or nonviable, and to compare it with the edge-based protocol of Go et al.~\cite{go2019} on images with and without lumped pollen. The details of this approach are described in Section~\ref{ch:methodology}.
